% GENERATED FILE. Edit canonical lessons, then run npm run generate:books.
% canonical-source-hash: fnv1a64:5ac1ec4206386641
% canonical-lessons: TA-C339-maviran, TA-C339-viram, TA-C339-tiyaki, TA-C339-munnoti, TA-C339-natukal, TA-R339-review-333-339

\chapter{Heroes Remembered}
\label{ch:ta-heroes-remembered}
\hlchaptermodality{339}

\begin{chapteropening}
\textbf{By the end of this chapter:} \emph{I can talk about heroes and how they are remembered.}

Complete the last lesson of chapter 339: I can talk about heroes and how they are remembered.
\end{chapteropening}

\section[\texorpdfstring{\ta{மாவீரன்} (māvīraṉ)}{மாவீரன் (māvīraṉ)}]{\ta{மாவீரன்} (māvīraṉ) — a great hero}

\label{lesson:TA-C339-maviran}



\begin{quote}
Before the new one: say the Tamil for a revolt, then the Tamil for to jail.
\end{quote}

\subsection*{You'll want to know: \ta{மாவீரன்}}
\textbf{\ta{மாவீரன்}} — \emph{māvīraṉ} — \textquotedblleft{}a great hero\textquotedblright{}.

\begin{quote}
\ta{மக்கள்} \ta{இன்றும்} \ta{அந்த} \ta{மாவீரனை} \ta{நினைக்கிறார்கள்}.
\end{quote}

\begin{quote}
\emph{makkaḷ iṉṟum anta māvīraṉai niṉaikkiṟārkaḷ.}
\end{quote}

\begin{quote}
People still remember that great hero today.
\end{quote}

\ta{மா}, great + \ta{வீரன்}, hero.

\begin{cousinweb}
\textbf{\ta{மாவீரன்}} (\emph{māvīraṉ}) is built from \ta{மா}, \textquotedblleft{}great\textquotedblright{}, and \ta{வீரன்}, \textquotedblleft{}hero\textquotedblright{}, borrowed from Sanskrit \emph{vīra}, \textquotedblleft{}man, hero\textquotedblright{} (CDIAL 12056).
\end{cousinweb}

\begin{grammarlens}[title={What you've built}]
One more word for telling what happened.
\end{grammarlens}

\subsection*{Your turn}
\begin{itemize}
\raggedright
  \item \emph{Say it:} \emph{māvīraṉ}
  \item \emph{Say it:} \emph{māvīraṉ}, once more
  \item \emph{Say it:} \emph{makkaḷ iṉṟum anta māvīraṉai niṉaikkiṟārkaḷ.} — People still remember that great hero today.
  \item \emph{Recall:} say the Tamil for a revolt, then the Tamil for to jail, then say \emph{māvīraṉ} again
\end{itemize}

\begin{tcolorbox}[breakable,colback=teal!4,colframe=teal!35!black,title={Before you move on}]
What does \ta{மாவீரன்} mean? (\textquotedblleft{}A great hero\textquotedblright{}.) Say it once more.
\end{tcolorbox}
\section[\texorpdfstring{\ta{வீரம்} (vīram)}{வீரம் (vīram)}]{\ta{வீரம்} (vīram) — bravery; valour}

\label{lesson:TA-C339-viram}



\begin{quote}
Before the new one: say the Tamil for to jail, then the Tamil for a great hero.
\end{quote}

\subsection*{You'll want to know: \ta{வீரம்}}
\textbf{\ta{வீரம்}} — \emph{vīram} — \textquotedblleft{}bravery; valour\textquotedblright{}.

\begin{quote}
\ta{அவருடைய} \ta{வீரத்தைப்} \ta{பற்றிப்} \ta{பல} \ta{பாடல்கள்} \ta{இருக்கின்றன}.
\end{quote}

\begin{quote}
\emph{avaruṭaiya vīrattaip paṟṟip pala pāṭalkaḷ irukkiṉṟaṉa.}
\end{quote}

\begin{quote}
There are many songs about his bravery.
\end{quote}

A -\ta{அம்} noun: \ta{வீரத்தை}, \ta{வீரத்தின்}.

\begin{cousinweb}
\textbf{\ta{வீரம்}} (\emph{vīram}) is built on Sanskrit \emph{vīra}, \textquotedblleft{}man, hero\textquotedblright{} (CDIAL 12056), with the Tamil -\ta{அம்} of nouns.
\end{cousinweb}

\begin{grammarlens}[title={What you've built}]
One more word for telling what happened.
\end{grammarlens}

\subsection*{Your turn}
\begin{itemize}
\raggedright
  \item \emph{Say it:} \emph{vīram}
  \item \emph{Say it:} \emph{vīram}, once more
  \item \emph{Say it:} \emph{avaruṭaiya vīrattaip paṟṟip pala pāṭalkaḷ irukkiṉṟaṉa.} — There are many songs about his bravery.
  \item \emph{Recall:} say the Tamil for to jail, then the Tamil for a great hero, then say \emph{vīram} again
\end{itemize}

\begin{tcolorbox}[breakable,colback=teal!4,colframe=teal!35!black,title={Before you move on}]
What does \ta{வீரம்} mean? (\textquotedblleft{}Bravery; valour\textquotedblright{}.) Say it once more.
\end{tcolorbox}
\section[\texorpdfstring{\ta{தியாகி} (tiyāki)}{தியாகி (tiyāki)}]{\ta{தியாகி} (tiyāki) — a martyr; one who sacrifices}

\label{lesson:TA-C339-tiyaki}



\begin{quote}
Before the new one: say the Tamil for a great hero, then the Tamil for bravery.
\end{quote}

\subsection*{You'll want to know: \ta{தியாகி}}
\textbf{\ta{தியாகி}} — \emph{tiyāki} — \textquotedblleft{}a martyr; one who sacrifices\textquotedblright{}.

\begin{quote}
\ta{நாட்டுக்காக} \ta{உயிரைக்} \ta{கொடுத்த} \ta{தியாகிகளை} \ta{நினைவுகூர்ந்தோம்}.
\end{quote}

\begin{quote}
\emph{nāṭṭukkāka uyiraik koṭutta tiyākikaḷai niṉaivukūrntōm.}
\end{quote}

\begin{quote}
We remembered the martyrs who gave their lives for the country.
\end{quote}

One who gives up everything for a cause.

\begin{cousinweb}
\textbf{\ta{தியாகி}} (\emph{tiyāki}) is borrowed from Sanskrit \emph{tyāgin}, \textquotedblleft{}one who abandons, gives up\textquotedblright{} (Monier-Williams).
\end{cousinweb}

\begin{grammarlens}[title={What you've built}]
One more word for telling what happened.
\end{grammarlens}

\subsection*{Your turn}
\begin{itemize}
\raggedright
  \item \emph{Say it:} \emph{tiyāki}
  \item \emph{Say it:} \emph{tiyāki}, once more
  \item \emph{Say it:} \emph{nāṭṭukkāka uyiraik koṭutta tiyākikaḷai niṉaivukūrntōm.} — We remembered the martyrs who gave their lives for the country.
  \item \emph{Recall:} say the Tamil for a great hero, then the Tamil for bravery, then say \emph{tiyāki} again
\end{itemize}

\begin{tcolorbox}[breakable,colback=teal!4,colframe=teal!35!black,title={Before you move on}]
What does \ta{தியாகி} mean? (\textquotedblleft{}A martyr; one who sacrifices\textquotedblright{}.) Say it once more.
\end{tcolorbox}
\section[\texorpdfstring{\ta{முன்னோடி} (muṉṉōṭi)}{முன்னோடி (muṉṉōṭi)}]{\ta{முன்னோடி} (muṉṉōṭi) — a pioneer; a forerunner}

\label{lesson:TA-C339-munnoti}



\begin{quote}
Before the new one: say the Tamil for bravery, then the Tamil for a martyr.
\end{quote}

\subsection*{You'll want to know: \ta{முன்னோடி}}
\textbf{\ta{முன்னோடி}} — \emph{muṉṉōṭi} — \textquotedblleft{}a pioneer; a forerunner\textquotedblright{}.

\begin{quote}
\ta{அவர்} \ta{தமிழ்} \ta{நாடகத்தின்} \ta{முன்னோடி}.
\end{quote}

\begin{quote}
\emph{avar tamiḻ nāṭakattiṉ muṉṉōṭi.}
\end{quote}

\begin{quote}
He is a pioneer of Tamil drama.
\end{quote}

\ta{முன்}, front + \ta{ஓடி}, one who runs.

\begin{cousinweb}
\textbf{\ta{முன்னோடி}} (\emph{muṉṉōṭi}) is a Tamil compound built from \ta{முன்}, \textquotedblleft{}in front\textquotedblright{} (DEDR 5020), and \ta{ஓடி}, \textquotedblleft{}one who runs\textquotedblright{}, from \ta{ஓடு} (DEDR 1041).
\end{cousinweb}

\begin{grammarlens}[title={What you've built}]
One more word for telling what happened.
\end{grammarlens}

\subsection*{Your turn}
\begin{itemize}
\raggedright
  \item \emph{Say it:} \emph{muṉṉōṭi}
  \item \emph{Say it:} \emph{muṉṉōṭi}, once more
  \item \emph{Say it:} \emph{avar tamiḻ nāṭakattiṉ muṉṉōṭi.} — He is a pioneer of Tamil drama.
  \item \emph{Recall:} say the Tamil for bravery, then the Tamil for a martyr, then say \emph{muṉṉōṭi} again
\end{itemize}

\begin{tcolorbox}[breakable,colback=teal!4,colframe=teal!35!black,title={Before you move on}]
What does \ta{முன்னோடி} mean? (\textquotedblleft{}A pioneer; a forerunner\textquotedblright{}.) Say it once more.
\end{tcolorbox}
\section[\texorpdfstring{\ta{நடுகல்} (naṭukal)}{நடுகல் (naṭukal)}]{\ta{நடுகல்} (naṭukal) — a hero stone}

\label{lesson:TA-C339-natukal}



\begin{quote}
Before the new one: say the Tamil for a martyr, then the Tamil for a pioneer.
\end{quote}

\subsection*{You'll want to know: \ta{நடுகல்}}
\textbf{\ta{நடுகல்}} — \emph{naṭukal} — \textquotedblleft{}a hero stone\textquotedblright{}.

\begin{quote}
\ta{ஊருக்கு} \ta{வெளியே} \ta{ஒரு} \ta{பழைய} \ta{நடுகல்} \ta{இருக்கிறது}.
\end{quote}

\begin{quote}
\emph{ūrukku veḷiyē oru paḻaiya naṭukal irukkiṟatu.}
\end{quote}

\begin{quote}
There is an old hero stone outside the village.
\end{quote}

A stone set up for a fallen hero, often carved with his figure.

\begin{cousinweb}
\textbf{\ta{நடுகல்}} (\emph{naṭukal}) is a Tamil compound built from \ta{நடு}, \textquotedblleft{}to set up, plant\textquotedblright{} (DEDR 3583), and \ta{கல்}, \textquotedblleft{}stone\textquotedblright{} (DEDR 1298).
\end{cousinweb}

\begin{grammarlens}[title={What you've built}]
One more word for telling what happened.
\end{grammarlens}

\subsection*{Your turn}
\begin{itemize}
\raggedright
  \item \emph{Say it:} \emph{naṭukal}
  \item \emph{Say it:} \emph{naṭukal}, once more
  \item \emph{Say it:} \emph{ūrukku veḷiyē oru paḻaiya naṭukal irukkiṟatu.} — There is an old hero stone outside the village.
  \item \emph{Recall:} say the Tamil for a martyr, then the Tamil for a pioneer, then say \emph{naṭukal} again
\end{itemize}

\begin{tcolorbox}[breakable,colback=teal!4,colframe=teal!35!black,title={Before you move on}]
What does \ta{நடுகல்} mean? (\textquotedblleft{}A hero stone\textquotedblright{}.) Say it once more.
\end{tcolorbox}
\section[(dialogue)]{Looking back: from The Battlefield to Heroes Remembered}

\label{lesson:TA-R339-review-333-339}



\begin{quote}
Say the Tamil for a battlefield and a soldier.
\end{quote}

\subsection*{Your turn}
Say these aloud:

\begin{itemize}
\raggedright
  \item a battlefield, a soldier, a sword, a shield, armour
  \item an arrow, a chariot, a horseman, a war drum, a commander
  \item to besiege, a moat, a tunnel, a fortification, a spy
  \item to bring down, to capture, to surrender, to submit, to occupy
  \item a conspiracy, a scheme, to betray, to deceive, a disguise
  \item a slave, liberation, a revolution, a revolt, to jail
  \item a great hero, bravery, a martyr, a pioneer, a hero stone
\end{itemize}

\begin{tcolorbox}[breakable,colback=teal!4,colframe=teal!35!black,title={Before you move on}]
A battlefield? (\textbf{\ta{போர்க்களம்}}, \emph{pōrkkaḷam}.) A soldier? (\textbf{\ta{போர்வீரர்}}, \emph{pōrvīrar}.) A sword? (\textbf{\ta{வாள்}}, \emph{vāḷ}.) A shield? (\textbf{\ta{கேடயம்}}, \emph{kēṭayam}.) Armour? (\textbf{\ta{கவசம்}}, \emph{kavacam}.) An arrow? (\textbf{\ta{அம்பு}}, \emph{ampu}.) A chariot? (\textbf{\ta{தேர்}}, \emph{tēr}.) A horseman? (\textbf{\ta{குதிரைவீரன்}}, \emph{kutiraivīraṉ}.) A war drum? (\textbf{\ta{முரசு}}, \emph{muracu}.) A commander? (\textbf{\ta{தளபதி}}, \emph{taḷapati}.) To besiege? (\textbf{\ta{முற்றுகையிடு}}, \emph{muṟṟukaiyiṭu}.) A moat? (\textbf{\ta{அகழி}}, \emph{akaḻi}.) A tunnel? (\textbf{\ta{சுரங்கம்}}, \emph{curaṅkam}.) A fortification? (\textbf{\ta{அரண்}}, \emph{araṇ}.) A spy? (\textbf{\ta{ஒற்றன்}}, \emph{oṟṟaṉ}.) To bring down? (\textbf{\ta{வீழ்த்து}}, \emph{vīḻttu}.) To capture? (\textbf{\ta{கைப்பற்று}}, \emph{kaippaṟṟu}.) To surrender? (\textbf{\ta{சரணடை}}, \emph{caraṇaṭai}.) To submit? (\textbf{\ta{அடிபணி}}, \emph{aṭipaṇi}.) To occupy? (\textbf{\ta{ஆக்கிரமி}}, \emph{ākkirami}.) A conspiracy? (\textbf{\ta{சதி}}, \emph{cati}.) A scheme? (\textbf{\ta{சூழ்ச்சி}}, \emph{cūḻcci}.) To betray? (\textbf{\ta{காட்டிக்கொடு}}, \emph{kāṭṭikkoṭu}.) To deceive? (\textbf{\ta{வஞ்சி}}, \emph{vañci}.) A disguise? (\textbf{\ta{மாறுவேடம்}}, \emph{māṟuvēṭam}.) A slave? (\textbf{\ta{அடிமை}}, \emph{aṭimai}.) Liberation? (\textbf{\ta{விடுதலை}}, \emph{viṭutalai}.) A revolution? (\textbf{\ta{புரட்சி}}, \emph{puraṭci}.) A revolt? (\textbf{\ta{கலகம்}}, \emph{kalakam}.) To jail? (\textbf{\ta{சிறையிலடை}}, \emph{ciṟaiyilaṭai}.) A great hero? (\textbf{\ta{மாவீரன்}}, \emph{māvīraṉ}.) Bravery? (\textbf{\ta{வீரம்}}, \emph{vīram}.) A martyr? (\textbf{\ta{தியாகி}}, \emph{tiyāki}.) A pioneer? (\textbf{\ta{முன்னோடி}}, \emph{muṉṉōṭi}.) A hero stone? (\textbf{\ta{நடுகல்}}, \emph{naṭukal}.)
\end{tcolorbox}
